\documentclass[reqno]{amsart}
\usepackage[T1]{fontenc}
\usepackage{amsmath,amssymb,amsthm}
\usepackage{mathtools}
\usepackage{natbib}

\usepackage{hyperref}
\hypersetup{
    colorlinks=true,       % False for boxes; true for colored text
    linkcolor=blue,        % Color for internal links
    filecolor=magenta,     % Color for file links
    urlcolor=cyan,         % Color for external links
    citecolor=green,       % Color for citations
}

\newtheorem{theorem}{Theorem}[section]
\newtheorem{proposition}[theorem]{Proposition}
\newtheorem{lemma}[theorem]{Lemma}
\newtheorem{corollary}[theorem]{Corollary}
\theoremstyle{definition}
\newtheorem{definition}[theorem]{Definition}
\newtheorem{remark}[theorem]{Remark}
\newtheorem{question}{Question}

\newcommand{\norm}[1]{\left\lVert #1 \right\rVert}
\newcommand{\abs}[1]{\left\lvert #1 \right\rvert}
\newcommand{\ip}[2]{\langle #1, #2 \rangle}
\newcommand{\ext}{\operatorname{ext}}
\newcommand{\co}{\operatorname{co}}
\newcommand{\R}{\mathbb{R}}
\newcommand{\inner}[2]{\langle #1,#2\rangle}

\title{Convexity and Boundary Positivity of the Hessian}
\author{Nam Le}
\address{Department of Computer Science,
    Faculty of Information Technology, University of Science, Vietnam National University,
    227 Nguyen Van Cu Street, Cho Quan Ward, Ho Chi Minh City}
\email{lnnam@hcmus.edu.vn}
\date{\today}

\begin{document}

\maketitle

\begin{abstract}
Let $C \subset \R^N$ be a closed convex set with nonempty interior, and let
$f \colon C \to \R$ be convex and twice differentiable on $C$. It is standard
that $\nabla^2 f(x) \succeq 0$ for every $x \in \operatorname{int}(C)$, since the
usual open-set characterization of convexity by the positivity of the Hessian
applies on the interior. The main issue is whether the same conclusion must
remain true at boundary points when no continuity of the Hessian is assumed up
to $\partial C$. This note sets up that question precisely, records the standard
facts that are already available, and isolates the geometric distinction between
positivity in all ambient directions and positivity only along feasible boundary
directions.
\end{abstract}

\section{Introduction}

For convex functions defined on open subsets of $\R^N$, second-order
information is well understood: if $f$ is twice differentiable and convex on an
open set $U$, then the Hessian satisfies
\begin{equation}
    \nabla^2 f(x) \succeq 0 \qquad \forall x \in U.
\end{equation}
In the present note, the domain is not assumed open. Instead, we consider a
closed convex set $C \subset \R^N$ with nonempty interior and a function
$f \colon C \to \R$ that is convex on $C$.

The interior of the problem is immediate: restricting attention to
$\operatorname{int}(C)$ gives the standard conclusion
\begin{equation}
    \nabla^2 f(x) \succeq 0 \qquad \forall x \in \operatorname{int}(C).
\end{equation}
The difficulty lies at boundary points. If one strengthens the regularity
assumption and requires that $f$ be of class $C^2$ on $C$, then the boundary
question is again straightforward, because the Hessian can be obtained as a
limit of interior Hessians. Without continuity of $\nabla^2 f$ up to the
boundary, however, the situation is less clear.

This leads to two closely related questions. The first is whether convexity and
twice differentiability on a closed convex set already force the Hessian to be
positive semidefinite at every boundary point. The second is whether a twice
differentiable convex function must automatically have continuous Hessian, which
would reduce the first question to the familiar interior argument. The purpose
of the present draft is not yet to settle these questions, but to formulate them
precisely and to set up the notation and background needed for a later analysis.

There is also a geometric reason to separate these issues carefully. At a point
$x \in \partial C$, convexity only tests the behavior of $f$ along directions
that keep one inside the set. Consequently, even if second-order information is
available at $x$, one must distinguish between positivity of the quadratic form
$h \mapsto \inner{\nabla^2 f(x) h}{h}$ on all of $\R^N$ and positivity only on
the cone of feasible directions determined by $C$ at $x$. This distinction is
one of the central themes in what follows.

\section{Background and Preliminaries}

We fix a closed convex set $C \subset \R^N$ with $\operatorname{int}(C)
\neq \varnothing$ and a function $f \colon C \to \R$.

\subsection{Convexity on a closed convex set}

The function $f$ is said to be convex on $C$ if
\begin{equation}
    f((1-\alpha)x+\alpha y) \leq (1-\alpha)f(x)+\alpha f(y)
\end{equation}
for every $x,y \in C$ and every $\alpha \in [0,1]$. Since $C$ is convex, each
segment $[x,y]$ is contained in $C$, so the restriction of $f$ to any such
segment is a one-variable convex function.

Two standard facts will be used repeatedly.

\begin{proposition}
If $f$ is convex on $C$, then $f$ is continuous on $\operatorname{int}(C)$.
Moreover, if $f$ is differentiable on $\operatorname{int}(C)$, then
$\nabla f$ is continuous on $\operatorname{int}(C)$.
\end{proposition}

The second statement is a standard consequence of the monotonicity of the
subdifferential and the fact that, for a differentiable convex function, the
subdifferential reduces to the singleton $\{\nabla f(x)\}$.

\subsection{Twice differentiability on $C$}

Because the domain is not open, the notion of twice differentiability at a
boundary point must be stated explicitly. In this note, we say that $f$ is
twice differentiable at $x \in C$ if there exist a vector $\nabla f(x) \in \R^N$
and a symmetric matrix $\nabla^2 f(x) \in \R^{N \times N}$ such that
\begin{equation}
    f(x+h)=f(x)+\inner{\nabla f(x)}{h}+\tfrac12 \inner{\nabla^2 f(x)h}{h}+o(\|h\|^2)
\end{equation}
as $h \to 0$ under the constraint that $x+h \in C$. Equivalently, the remainder
is required to satisfy
\begin{equation}
    f(x+h)-f(x)-\inner{\nabla f(x)}{h}-\tfrac12 \inner{\nabla^2 f(x)h}{h}=o(\|h\|^2)
\end{equation}
for admissible increments $h$ with $x+h \in C$.

This is a relative notion of second-order differentiability along the set $C$.
When $x \in \operatorname{int}(C)$, it coincides with the usual definition on an
open neighborhood of $x$. At a boundary point, it captures second-order
expansions only along directions that remain in the domain.

\subsection{Feasible directions and tangent cones}

For $x \in C$, the feasible-direction cone of $C$ at $x$ is
\begin{equation}
    T_C(x):=\overline{\{t(y-x): t \geq 0,\ y \in C\}}.
\end{equation}
Equivalently, $v \in T_C(x)$ if there exist sequences $t_k \downarrow 0$ and
$v_k \to v$ such that $x+t_k v_k \in C$ for all $k$. When $x \in
\operatorname{int}(C)$, one has $T_C(x)=\R^N$; at a boundary point, this cone
encodes the directions in which one can move while remaining in $C$.

This cone is relevant because convexity of $f$ on $C$ can only be tested along
segments contained in $C$. Thus, at $x \in \partial C$, any immediate
second-order consequence of convexity is expected to involve the quadratic form
$h \mapsto \inner{\nabla^2 f(x)h}{h}$ at least on $T_C(x)$, whereas positivity in
arbitrary ambient directions is a stronger requirement.

\subsection{Standard second-order criterion on the interior}

The following is the classical open-set result applied to the interior of $C$.

\begin{proposition}
Suppose that $f$ is convex on $C$ and twice differentiable on
$\operatorname{int}(C)$. Then
\begin{equation}
    \nabla^2 f(x) \succeq 0 \qquad \forall x \in \operatorname{int}(C).
\end{equation}
\end{proposition}

Indeed, each point of $\operatorname{int}(C)$ has a neighborhood contained in
$C$, so the standard theorem for convex $C^2$ functions on open sets applies.

\begin{remark}
If one additionally assumes that $f$ is of class $C^2$ on the whole of $C$ in
the sense that the Hessian extends continuously to boundary points, then the
interior positivity immediately passes to every $x \in C$ by taking limits from
$\operatorname{int}(C)$.
\end{remark}

\section{Problem Statement}

We can now formulate the main issue precisely.

\begin{question}[Boundary positivity of the Hessian]
Let $C \subset \R^N$ be a closed convex set with nonempty interior, and let
$f \colon C \to \R$ be convex and twice differentiable on $C$ in the relative
sense above. Must one have
\begin{equation}
    \nabla^2 f(x) \succeq 0 \qquad \forall x \in C?
\end{equation}
Equivalently, does convexity on $C$ force the Hessian at a boundary point to be
positive semidefinite on all ambient directions of $\R^N$?
\end{question}

This question is nontrivial only on $\partial C$, since the conclusion is already
known on $\operatorname{int}(C)$. It is also already settled under the stronger
assumption that $f$ is $C^2$ on $C$.

\begin{question}[Automatic continuity of the Hessian]
If $f$ is convex and twice differentiable on $C$, must the mapping
$x \mapsto \nabla^2 f(x)$ be continuous, at least on compact subsets of
$\operatorname{int}(C)$ or up to the boundary under the relative notion of
twice differentiability?
\end{question}

An affirmative answer to the second question would give a direct positive answer
to the first one, but it is not clear a priori that convexity plus existence of
pointwise second derivatives should force such continuity.

There is also a weaker boundary formulation which is geometrically natural.

\begin{question}[Positivity on feasible directions]
For $x \in \partial C$, does convexity at least imply that
\begin{equation}
    \inner{\nabla^2 f(x)h}{h} \geq 0 \qquad \forall h \in T_C(x)?
\end{equation}
If so, to what extent can this feasible-direction positivity be upgraded to full
positive semidefiniteness on $\R^N$?
\end{question}

This reformulation isolates the geometric core of the problem. Convexity of the
restriction of $f$ to segments in $C$ naturally gives one-dimensional
second-order inequalities only in directions allowed by the set. The passage
from such restricted inequalities to an ambient matrix inequality is exactly the
point that must be clarified.

\begin{remark}[Expected route for the sequel]
A later section should examine whether the correct boundary statement is full
positivity of $\nabla^2 f(x)$, positivity only on $T_C(x)$, or some intermediate
condition depending on the local geometry of $\partial C$. It will also be useful
to separate smooth boundary points from corners or other nonsmooth points of the
set, since the tangent structure may behave differently in those cases.
\end{remark}

\section{Next Step Placeholder}

The next stage of the note should provide either a proof or a counterexample for
the boundary question, together with a separate discussion of the continuity of
the Hessian for twice differentiable convex functions.

\end{document}
